\question{7.m3}{
    Het aantal ongevallen per maand op een kruispunt kan worden beschreven door een Poissonverdeling met $\lambda=0.5$.
    De standaarddeviatie van het aantal ongelukken per jaar bedraagt daarom $\ldots$
}
\begin{enumerate}[label=(\alph*)]
    \item $6$.
    \item $\sqrt{6}$.
    \item $\frac{0.5}{\sqrt{6}}$.
    \item $0.5 \times \sqrt{12}$.
\end{enumerate}
\answer{
    Een jaar heeft $12$ maanden, dus we bekijken een interval van twaalf maanden met een tijdseenheid van een maand, oftewel $t=12$.
    Het aantal ongevallen per jaar kan in dat geval worden beschreven door een Poissonverdeling met $\mu = \lambda \cdot t = 0.5 \cdot 12 = 6$.
    De standaarddeviatie van een $\poisson{\mu}$-verdeling is gelijk aan $\sqrt{\mu}$, oftewel $\sqrt{\mu} = \sqrt{6}$.
    Het juiste antwoord is dus (b).
}